% Unit 045 terjemahan Bahasa Indonesia dari chapter3.tex baris 2714--2935.
% Sumber: Methods of Algebra, Volume 2, commit
% 9a5803ff77dd3257484cb177f851a73770a59dd3 (CC BY 4.0).
% stable-unit-id: o014.aljabr2.chapter3.ext-tor
% Cakupan: Bagian 3.14 lengkap.

% segment-id: o014.li2.chapter3.ext-tor.h001
\section{Contoh: \texorpdfstring{$\Ext$}{Ext} dan \texorpdfstring{$\Tor$}{Tor}}\label{sec:Ext-Tor}
\hypertarget{o014.aljabr2.chapter3.ext-tor}{ }

% segment-id: o014.li2.chapter3.ext-tor.p001
Kita akan mengkaji funktor turunan dari $\Hom$ dan $\otimes$. Sebaiknya kita
memperkenalkan beberapa konsep bantu terlebih dahulu.

% segment-id: o014.li2.chapter3.ext-tor.de001
\begin{definisi}\label{def:balanced-functor}
\index{bifunktor!seimbang (balanced)}
Misalkan $\mathcal{A}_1$, $\mathcal{A}_2$, dan $\mathcal{B}$ merupakan
kategori abelian, sedangkan bifunktor $F: \mathcal{A}_1 \times \mathcal{A}_2
\to \mathcal{B}$ (Konvensi \ref{con:bifunctor}) merupakan funktor eksak kiri
terhadap kedua variabel. Jika untuk setiap objek injektif $I_i \in
\Obj(\mathcal{A}_i)$ (dengan $i=1,2$),
\[
F(I_1, \cdot): \mathcal{A}_2 \to \mathcal{B} \quad\text{dan}\quad
F(\cdot, I_2): \mathcal{A}_1 \to \mathcal{B}
\]
keduanya merupakan funktor eksak, maka $F$ disebut \emph{seimbang}.

Secara dual, jika $F$ eksak kanan terhadap kedua variabel, dan untuk setiap
objek projektif $P_i \in \Obj(\mathcal{A}_i)$ funktor $F(P_1,\cdot)$ dan
$F(\cdot,P_2)$ keduanya eksak, maka $F$ juga disebut seimbang.
\end{definisi}

% segment-id: o014.li2.chapter3.ext-tor.p002
Berdasarkan dualitas, teorema berikut hanya dinyatakan untuk bifunktor yang
eksak kiri terhadap kedua variabel.

% segment-id: o014.li2.chapter3.ext-tor.th001
\begin{teorema}\label{prop:balanced-primer}
\index[sym1]{RInF@$\mathrm{R}_{\mathrm{I}}^n F$, $\mathrm{R}_{\mathrm{II}}^n F$}
Misalkan $\mathcal{A}_1$, $\mathcal{A}_2$, dan $\mathcal{B}$ merupakan
kategori abelian; $\mathcal{A}_1$ dan $\mathcal{A}_2$ mempunyai cukup banyak
objek injektif. Misalkan bifunktor $F: \mathcal{A}_1\times\mathcal{A}_2\to
\mathcal{B}$ eksak kiri terhadap kedua variabel. Dengan demikian, kita dapat
mengambil funktor turunan kanan terhadap setiap variabel. Untuk $X_i \in
\Obj(\mathcal{A}_i)$ yang tetap, $i=1,2$, tuliskan
\[
\mathrm{R}_{\mathrm{I}}^n F(\cdot,X_2):\mathcal{A}_1\to\mathcal{B},\qquad
\mathrm{R}_{\mathrm{II}}^n F(X_1,\cdot):\mathcal{A}_2\to\mathcal{B}.
\]
Jika $F$ seimbang, maka terdapat isomorfisme kanonik bifunktor
\[
\mathrm{R}_{\mathrm{I}}^nF(X_1,X_2)\simeq
\mathrm{R}_{\mathrm{II}}^nF(X_1,X_2).
\]
\end{teorema}

% segment-id: o014.li2.chapter3.ext-tor.pf001
\begin{bukti}
Argumen berikut bergantung pada teori bikompleks; lihat
\S\ref{sec:double-cplx} dan \S\ref{sec:double-cplx-coh}. Untuk $i=1,2$,
ambil resolusi injektif $0\to X_i\to I_i^0\to I_i^1\to\cdots$; untuk
$n<0$ tetapkan $I_i^n:=0$. Ini mendefinisikan bikompleks
$F\left(I_1^\bullet,I_2^\bullet\right)$.

Di sisi lain, definisikan kompleks $F(X_1,I_2^\bullet)$ dan
$F(I_1^\bullet,X_2)$. Pandang keduanya sebagai bikompleks yang masing-masing
terkonsentrasi pada sumbu vertikal $(0,\bullet)$ dan sumbu horizontal
$(\bullet,0)$. Karena
$X_i\rightiso\Ker[I_i^0\to I_i^1]\hookrightarrow I_i^0$, diperoleh dalam
$\cate{C}^2_f(\mathcal{B})$ morfisme
\[
F(X_1,I_2^\bullet)\xrightarrow{\lambda}
F\left(I_1^\bullet,I_2^\bullet\right)\xleftarrow{\rho}
F(I_1^\bullet,X_2).
\]
Kita akan menunjukkan bahwa $\tot(\lambda)$ dan $\tot(\rho)$ merupakan
kuasi-isomorfisme.

Pertama-tama, pertimbangkan $\lambda$. Ambil kohomologi horizontal
$\Hm_{\mathrm{I}}$ pada kedua sisinya; hal ini tidak mengubah
$F(X_1,I_2^\bullet)$ yang terletak pada sumbu vertikal. Karena
$F(\cdot,I_2^q)$ merupakan funktor eksak untuk setiap $q$, kita memperoleh
\[
\left(\Hm_{\mathrm{I}}F\left(I_1^\bullet,I_2^\bullet\right)\right)^{p,q}
=\begin{cases}
0,&p\neq0,\\
\Ker\left[F(I_1^0,I_2^q)\to F(I_1^1,I_2^q)\right]
=F(X_1,I_2^q),&p=0.
\end{cases}
\]
Jadi $\Hm_{\mathrm{I}}(\lambda)$ merupakan isomorfisme dari
$\Hm_{\mathrm{I}}F(X_1,I_2^\bullet)$ ke
$\Hm_{\mathrm{I}}F(I_1^\bullet,I_2^\bullet)$. Maka
$\Hm_{\mathrm{II}}\Hm_{\mathrm{I}}(\lambda)$ juga isomorfisme. Dengan
menerapkan Teorema \ref{prop:double-cplx-tot}, diperoleh bahwa
$\tot(\lambda)$ merupakan kuasi-isomorfisme.

Untuk $\rho$, argumen yang sama memberikan bahwa
$\Hm_{\mathrm{I}}\Hm_{\mathrm{II}}(\rho)$ merupakan isomorfisme; Teorema
\ref{prop:double-cplx-tot} lalu menyiratkan bahwa $\tot(\rho)$ juga merupakan
kuasi-isomorfisme.

Perhatikan bahwa kompleks total $F(X_1,I_2^\bullet)$ tetap merupakan
kompleks itu sendiri; mengambil $\Hm^n$ menghasilkan
$\mathrm{R}_{\mathrm{II}}^nF(X_1,X_2)$. Demikian pula, mengambil $\Hm^n$ dari
$F(I_1^\bullet,X_2)$ atau kompleks totalnya menghasilkan
$\mathrm{R}_{\mathrm{I}}^nF(X_1,X_2)$. Melalui kompleks total
$\tot\left(F(I_1^\bullet,I_2^\bullet)\right)$ sebagai perantara, diperoleh
\[
\mathrm{R}_{\mathrm{I}}^nF(X_1,X_2)\simeq
\mathrm{R}_{\mathrm{II}}^nF(X_1,X_2).
\]
Karena setiap morfisme $X_i\to Y_i$ dapat diangkat menjadi morfisme di antara
resolusi injektif, dan pengangkatan tersebut unik hingga homotopi, yang kita
dapatkan sebenarnya adalah isomorfisme bifunktor
$\mathrm{R}_{\mathrm{I}}^nF\simeq\mathrm{R}_{\mathrm{II}}^nF$.
\end{bukti}

% segment-id: o014.li2.chapter3.ext-tor.rem001
\begin{catatan}\label{rem:balanced-primer}
Argumen tersebut sesungguhnya memberikan teknik untuk meresolusikan kedua
variabel sekaligus, lalu ``menurunkan'' $F$ secara seimbang melalui kompleks
total dari bikompleks $F(I_1^\bullet,I_2^\bullet)$.
\end{catatan}

% segment-id: o014.li2.chapter3.ext-tor.p003
Sekarang ambil $\mathcal{A}$ sebagai kategori abelian dan pertimbangkan
bifunktor $\Hom=\Hom_{\mathcal{A}}:\mathcal{A}^{\opp}\times\mathcal{A}
\to\cate{Ab}$. Proposisi \ref{prop:Hom-left-exact} menyatakan bahwa
$\Hom$ eksak kiri terhadap kedua variabel. Untuk setiap $n\in\Z$ kita buat
definisi berikut.

% segment-id: o014.li2.chapter3.ext-tor.l001
\begin{itemize}
% segment-id: o014.li2.chapter3.ext-tor.i001
\item Jika $\mathcal{A}$ mempunyai cukup banyak objek projektif, definisikan
$\Ext^n_{\mathcal{A},\mathrm{I}}(X,Y):=\mathrm{R}^n_{\mathrm{I}}\Hom(X,Y)$,
yang ditentukan oleh resolusi projektif $X$ di $\mathcal{A}$.
% segment-id: o014.li2.chapter3.ext-tor.i002
\item Jika $\mathcal{A}$ mempunyai cukup banyak objek injektif, definisikan
$\Ext^n_{\mathcal{A},\mathrm{II}}(X,Y):=\mathrm{R}^n_{\mathrm{II}}\Hom(X,Y)$,
yang ditentukan oleh resolusi injektif $Y$ di $\mathcal{A}$.
\end{itemize}

% segment-id: o014.li2.chapter3.ext-tor.dp001
\begin{definisiproposisi}[funktor $\Ext$]\label{def:Ext-classical}
\index{Ext functor@funktor $\Ext$ ($\Ext$ functor)}
\index[sym1]{Extn@$\Ext^n$}
Bifunktor $\Hom_{\mathcal{A}}$ bersifat seimbang. Karena itu, selama kategori
abelian $\mathcal{A}$ mempunyai cukup banyak objek injektif dan projektif, untuk
setiap $n\in\Z$ terdapat bifunktor
\[
\Ext^n=\Ext^n_{\mathcal{A}}:\mathcal{A}^{\opp}\times\mathcal{A}\to\cate{Ab}
\]
dengan
$\Ext^n_{\mathcal{A}}(X,Y):=\Ext^n_{\mathrm{I}}(X,Y)\simeq
\Ext^n_{\mathrm{II}}(X,Y)$ (isomorfisme kanonik).
\end{definisiproposisi}

% segment-id: o014.li2.chapter3.ext-tor.pf002
\begin{bukti}
Keseimbangan tepat merupakan konsekuensi dari definisi objek injektif dan
projektif, sehingga Teorema \ref{prop:balanced-primer} dapat diterapkan.
\end{bukti}

% segment-id: o014.li2.chapter3.ext-tor.p004
Jika $R$ merupakan gelanggang dan $\mathcal{A}=R\dcate{Mod}$ atau
$\cated{Mod}R$, kita tulis $\Ext^n_{\mathcal{A}}$ sebagai $\Ext^n_R$.

% segment-id: o014.li2.chapter3.ext-tor.p005
Perhatikan bahwa $\Ext^0=\Hom$, sedangkan $(\Ext^n)_{n\geq0}$, masing-masing
terhadap variabel $X$ dan $Y$, mempunyai barisan eksak panjang berbentuk
% segment-id: o014.li2.chapter3.ext-tor.q001
\begin{gather*}
\cdots\to\Ext^{n-1}(X',Y)\xrightarrow{\delta^{n-1}}\Ext^n(X'',Y)\to
\Ext^n(X,Y)\to\Ext^n(X',Y)\to\cdots,\\
\cdots\to\Ext^{n-1}(X,Y'')\xrightarrow{\delta^{n-1}}\Ext^n(X,Y')\to
\Ext^n(X,Y)\to\Ext^n(X,Y'')\to\cdots,
\end{gather*}
di mana $0\to X'\to X\to X''\to0$ dan
$0\to Y'\to Y\to Y''\to0$ merupakan barisan eksak pendek yang diberikan.

% segment-id: o014.li2.chapter3.ext-tor.p006
Jika $\mathcal{A}$ merupakan kategori abelian $\Bbbk$-linear, nilai $\Ext^n$
tentu dapat ditingkatkan ke dalam kategori $\Bbbk\dcate{Mod}$.

% segment-id: o014.li2.chapter3.ext-tor.rem002
\begin{catatan}
Simbol $\Ext$ berarti ``ekstensi'', sebab elemen-elemen $\Ext^1(X,Y)$
berkorespondensi satu-satu dengan kelas isomorfisme ekstensi
$0\to Y\to E\to X\to0$ (yakni barisan eksak pendek); $\Ext^n$ yang lebih
tinggi mempunyai interpretasi serupa. Fakta ini sebaiknya dipahami melalui
kategori turunan; lihat \S\sourcecrossref{sec:Hom-Ext}{4.5}.
\end{catatan}

% segment-id: o014.li2.chapter3.ext-tor.p007
Berdasarkan karakterisasi keeksakan funktor melalui Korolari
\ref{prop:obstruction-exactness}, segera diperoleh

% segment-id: o014.li2.chapter3.ext-tor.l002
\begin{itemize}
% segment-id: o014.li2.chapter3.ext-tor.i003
\item objek $I$ injektif $\iff \Ext^1(\cdot,I)=0\iff
\Ext^{\geq1}(\cdot,I)=0$;
% segment-id: o014.li2.chapter3.ext-tor.i004
\item objek $P$ projektif $\iff \Ext^1(P,\cdot)=0\iff
\Ext^{\geq1}(P,\cdot)=0$.
\end{itemize}

% segment-id: o014.li2.chapter3.ext-tor.ex001
\begin{contoh}[$\HHm^n$ sebagai $\Ext^n$]\label{eg:HH-Ext}
\index[sym1]{HH}
Kembali ke kerangka dasar homologi dan kohomologi Hochschild pada
\S\ref{sec:HH}: $\Bbbk$ merupakan gelanggang komutatif dan $R$ merupakan
aljabar-$\Bbbk$; sekarang syaratkan pula bahwa $R$ projektif sebagai
modul-$\Bbbk$. Tuliskan $\otimes:=\otimes_{\Bbbk}$. Karena modul projektif
merupakan suku langsung dari modul bebas, $R^{\otimes n}$ tetap merupakan
modul-$\Bbbk$ projektif.
\footnote{Atau gunakan hubungan adjoin tensor \cite[Teorema 6.6.5]{Li1}.}
Jadikan $R\otimes R^{\otimes n}\otimes R$ sebagai modul-$R^e$ kiri, dengan
$R^e:=R\otimes R^{\opp}$. Karena
$R\otimes(\cdot)\otimes R:\Bbbk\dcate{Mod}\to R^e\dcate{Mod}$ merupakan
adjoin kiri dari funktor pelupa \cite[Korolari 6.6.8]{Li1}, funktor ini
membawa objek projektif ke objek projektif (Proposisi
\ref{prop:adjoint-injective-projective}). Maka Lema
\ref{prop:HH-bar-exactness} menunjukkan bahwa kompleks bar $\mathsf{B}R$
memberikan resolusi projektif $R$ sebagai modul-$R^e$ kiri,
\begin{equation}\label{eqn:bar-proj-resolution}
\cdots\to\mathsf{B}_1R\to\mathsf{B}_0R\to R\to0.
\end{equation}
Sekarang misalkan $M$ merupakan bimodul $(R,R)$. Dengan memasukkan definisi
awal kohomologi Hochschild \ref{def:HH}, diperoleh
\[
\HHm^n(M)=\Hm^n\left(\Hom_{R^e}(\mathsf{B}R,M)\right)
=\Ext^n_{R^e}(R,M).
\]
\end{contoh}

% segment-id: o014.li2.chapter3.ext-tor.p008
Topik berikutnya adalah funktor $\Tor$. Tetapkan gelanggang $R$. Kategori
modul $R\dcate{Mod}$ (atau $\cated{Mod}R$) mempunyai cukup banyak objek
projektif\footnote{Sesungguhnya, modul bebas selalu projektif; lihat
\cite[Contoh 6.9.6]{Li1}.}. Hasil kali tensor memberikan bifunktor
% segment-id: o014.li2.chapter3.ext-tor.q002
\[
\otimes_R:\cated{Mod}R\times R\dcate{Mod}\to\cate{Ab},
\]
yang eksak kanan terhadap setiap variabel (lihat \cite[Proposisi 6.9.2]{Li1}).
Terapkan versi eksak kanan dari teori sebelumnya untuk
mendefinisikan funktor turunan kiri terhadap kedua variabel,
$\mathrm{L}_{\mathrm{I},n}\otimes_R$ dan
$\mathrm{L}_{\mathrm{II},n}\otimes_R$; keduanya hanya tak nol untuk $n\geq0$.
Ingat juga definisi modul datar \cite[Definisi 6.9.4]{Li1}.

% segment-id: o014.li2.chapter3.ext-tor.dp002
\begin{definisiproposisi}[funktor $\Tor$]\label{def:Tor-classical}
\index{Tor functor@funktor $\Tor$ ($\Tor$ functor)}
\index[sym1]{Torn@$\Tor^R_n$}
Untuk setiap gelanggang $R$, bifunktor $\otimes_R$ bersifat seimbang. Maka,
untuk setiap $n\in\Z$ dapat didefinisikan bifunktor
\[
\Tor_n=\Tor_n^R:\cated{Mod}R\times R\dcate{Mod}\to\cate{Ab}
\]
dengan
$\Tor_n^R(X,Y):=(\mathrm{L}_{\mathrm{I},n}\otimes_R)(X,Y)\simeq
(\mathrm{L}_{\mathrm{II},n}\otimes_R)(X,Y)$ (isomorfisme kanonik).
\end{definisiproposisi}

% segment-id: o014.li2.chapter3.ext-tor.pf003
\begin{bukti}
Keseimbangan mengikuti dari fakta bahwa setiap modul projektif merupakan
modul datar \cite[Korolari 6.9.9]{Li1}. Terapkan versi funktor turunan kiri
dari Teorema \ref{prop:balanced-primer}.
\end{bukti}

% segment-id: o014.li2.chapter3.ext-tor.p009
Kita mempunyai $\Tor^R_0(X,Y)=X\dotimes{R}Y$, sedangkan $(\Tor_n)_{n\geq0}$
mempunyai barisan eksak panjang berikut:
% segment-id: o014.li2.chapter3.ext-tor.q003
\begin{gather*}
\cdots\to\Tor^R_{n+1}(X'',Y)\xrightarrow{\partial_{n+1}}\Tor^R_n(X',Y)\to
\Tor^R_n(X,Y)\to\Tor^R_n(X'',Y)\to\cdots,\\
\cdots\to\Tor^R_{n+1}(X,Y'')\xrightarrow{\partial_{n+1}}\Tor^R_n(X,Y')\to
\Tor^R_n(X,Y)\to\Tor^R_n(X,Y'')\to\cdots,
\end{gather*}
di mana $0\to X'\to X\to X''\to0$ dan
$0\to Y'\to Y\to Y''\to0$ merupakan barisan eksak pendek yang diberikan.

% segment-id: o014.li2.chapter3.ext-tor.p010
Kita mengulangi ungkapan yang serupa dengan kasus $\Ext$.

% segment-id: o014.li2.chapter3.ext-tor.l003
\begin{itemize}
% segment-id: o014.li2.chapter3.ext-tor.i005
\item Modul-$R$ kanan $M$ datar $\iff \Tor^R_1(M,\cdot)=0\iff
\Tor^R_{\geq1}(M,\cdot)=0$;
% segment-id: o014.li2.chapter3.ext-tor.i006
\item Modul-$R$ kiri $N$ datar $\iff \Tor^R_1(\cdot,N)=0\iff
\Tor^R_{\geq1}(\cdot,N)=0$.
\end{itemize}

% segment-id: o014.li2.chapter3.ext-tor.p011
Misalkan $M$ merupakan modul-$R$. Meniru definisi resolusi projektif,
\emph{resolusi datar}\index{resolusidatar@resolusi datar (flat resolution)}
$M$ didefinisikan sebagai barisan eksak berbentuk
% segment-id: o014.li2.chapter3.ext-tor.q004
\[
\cdots\to F^2\to F^1\to F^0\to M\to0,\qquad
\text{setiap $F^n$ merupakan modul-$R$ datar}.
\]

% segment-id: o014.li2.chapter3.ext-tor.p012
Dengan Korolari \ref{prop:acyclic-resolution} dan pengamatan di atas,
$\Tor^R_n(X,Y)$ juga dapat dihitung menggunakan resolusi datar $X$ atau $Y$.
Inilah teknik praktis untuk menghitung $\Tor$, sedangkan resolusi projektif
terutama berperan pada tingkat teoretis.

% segment-id: o014.li2.chapter3.ext-tor.p013
Perhitungan semacam ini secara alami juga dapat ``diseimbangkan''. Ambil
sembarang resolusi datar $\cdots\to P_1\to P_0\to X\to0$ dan
$\cdots\to Q_1\to Q_0\to Y\to0$. Bentuk bikompleks dalam arti kompleks rantai
$P_\bullet\dotimes{R}Q_\bullet$, dan tuliskan kompleks totalnya sebagai
$P\otimes Q$. Versi homologi dari Catatan \ref{rem:balanced-primer}
memberikan $\Tor_n^R(X,Y)\simeq\Hm_n(P\otimes Q)$; sifat yang diperlukan hanya
keeksakan $P_p\dotimes{R}(\cdot)$ dan $(\cdot)\dotimes{R}Q_q$, tepat sifat
yang dijamin oleh kedataran.

% segment-id: o014.li2.chapter3.ext-tor.rem003
\begin{catatan}[Struktur bimodul]\label{rem:ExtTor-bimodule}
Tetapkan gelanggang $A$ dan $B$; misalkan $X$ merupakan bimodul $(A,R)$ dan
$Y$ merupakan bimodul $(R,B)$. Karena $\Tor_n^R$ merupakan bifunktor,
$\Tor_n^R(X,Y)$ memiliki struktur bimodul kanonik $(A,B)$. Sebagai contoh,
pertimbangkan perkalian kiri oleh $A$: setiap $a\in A$ memberikan endomorfisme
$X$ sebagai modul-$R$ kanan, sehingga menginduksi endomorfisme
$L_a$ pada $\Tor_n^R(X,Y)$; berlaku $L_aL_{a'}=L_{aa'}$ dan
$L_1=\identity$. Perkalian kanan oleh $B$ diperlakukan serupa pada variabel
kedua. Untuk $n=0$, kita kembali pada struktur bimodul alami di
$X\dotimes{R}Y$.

Jika $Y$ dipilih bersama resolusi datarnya sebagai modul-$R$ kiri
$\cdots\to P_0\to Y\to0$, aksi perkalian kiri $A$ dapat langsung dibaca pada
kompleks $X\dotimes{R}P_\bullet$: aksi itu berasal dari perkalian kiri $A$
pada $X$. Perkalian kanan diperlakukan dengan cara yang sama.

Sebagai kasus khusus, jika dipilih homomorfisme gelanggang $R\to S$, maka
$\Tor_n^R(\cdot,S)$ (atau $\Tor_n^R(S,\cdot)$) meningkat menjadi funktor
$\cated{Mod}R\to\cated{Mod}S$ (atau $R\dcate{Mod}\to S\dcate{Mod}$).

Argumen serupa menunjukkan bahwa untuk $X$ yang merupakan bimodul $(R,A)$ dan
$Y$ yang merupakan bimodul $(R,B)$, setiap $\Ext_R^n(X,Y)$ secara alami
merupakan bimodul $(A,B)$.
\end{catatan}

% segment-id: o014.li2.chapter3.ext-tor.p014
Untuk gelanggang komutatif $R$, modul-modul atas $R$ tidak perlu dibedakan
menjadi kiri dan kanan; kategori yang bersesuaian ditulis $R\dcate{Mod}$. Maka
setiap modul-$R$ secara alami merupakan bimodul $(R,R)$, sehingga $\Tor_n^R$
menjadi bifunktor $(R\dcate{Mod})^2\to R\dcate{Mod}$.

% segment-id: o014.li2.chapter3.ext-tor.pr001
\begin{proposisi}\label{prop:Tor-symmetry}
Untuk gelanggang komutatif $R$, terdapat isomorfisme kanonik bifunktor
\[
\Tor_n^R(X,Y)\simeq\Tor_n^R(Y,X),
\]
untuk setiap $n\in\Z$.
\end{proposisi}

% segment-id: o014.li2.chapter3.ext-tor.pf004
\begin{bukti}
Ambil resolusi datar $P_\bullet$ dan $Q_\bullet$ masing-masing dari $X$ dan
$Y$. Tuliskan $P\otimes Q$ sebagai kompleks total dari bikompleks
$P_\bullet\dotimes{R}Q_\bullet$; definisikan $Q\otimes P$ secara serupa.
Kendala pertukaran hasil kali tensor \cite[Proposisi 6.5.14]{Li1} memberikan
isomorfisme kanonik
% segment-id: o014.li2.chapter3.ext-tor.q005
\[
P_p\dotimes{R}Q_q\simeq Q_q\dotimes{R}P_p,\qquad p,q\in\Z.
\]
Dengan kata lain,
$P_\bullet\dotimes{R}Q_\bullet\simeq
\mathrm{swap}(Q_\bullet\dotimes{R}P_\bullet)$. Terapkan Proposisi
\ref{prop:double-cplx-swap} untuk memperoleh isomorfisme kanonik kompleks
total $P\otimes Q\simeq Q\otimes P$. Ini membuktikan pernyataan.
\end{bukti}

% segment-id: o014.li2.chapter3.ext-tor.rem004
\begin{catatan}
Simbol $\Tor$ berarti ``torsi'', yang berasal dari contoh berikut. Ambil
ideal kiri $\mathfrak{a}$ pada gelanggang $R$, serta modul-$R$ kanan $X$
sembarang. Karena modul kiri bebas $R$ datar, terapkan Proposisi
\ref{prop:dimension-shifting} pada $0\to\mathfrak{a}\to R\to
R/\mathfrak{a}\to0$ untuk melakukan pergeseran dimensi. Diperoleh
% segment-id: o014.li2.chapter3.ext-tor.q006
\begin{align*}
\Tor^R_1(X,R/\mathfrak{a}) &\simeq \Ker\left[
X\dotimes{R}\mathfrak{a}\to X\dotimes{R}R\simeq X\right]\\
&\simeq \left\{\sum_i x_i\otimes a_i\in X\dotimes{R}\mathfrak{a}:
\sum_i x_i a_i=0\right\},\\
\Tor^R_n(X,R/\mathfrak{a})&\simeq\Tor^R_{n-1}(X,\mathfrak{a}),\qquad(n>1).
\end{align*}
Jika selanjutnya $t\in R$ bukan pembagi nol kanan, dan diambil
$\mathfrak{a}:=Rt\leftiso R$ (memetakan $r$ ke $rt$), maka
$\Tor^R_1(X,R/Rt)\simeq\{x\in X:xt=0\}$, yakni submodul torsi-$t$ dari $X$;
untuk $n>1$, $\Tor^R_n(X,R/Rt)=0$.
\end{catatan}

% segment-id: o014.li2.chapter3.ext-tor.ex002
\begin{contoh}[$\HHm_n$ sebagai $\Tor_n$]\label{eg:HH-Tor}
\index[sym1]{HH}
Lanjutkan langkah awal Contoh \ref{eg:HH-Ext}, tetapi lemahkan syarat pada
$R$ menjadi: $R$ merupakan modul datar atas $\Bbbk$. Maka $R^{\otimes n}$
juga datar. Kendala asosiatif
$(\cdot)\dotimes{R^e}(R\otimes R^{\otimes n}\otimes R)\simeq
(\cdot)\otimes R^{\otimes n}$ menunjukkan bahwa
$R\otimes R^{\otimes n}\otimes R$ merupakan modul-$R^e$ kiri yang datar.
Jadi, $\mathsf{B}R$ merupakan resolusi datar $R$ sebagai modul-$R^e$ kiri,
seperti pada \eqref{eqn:bar-proj-resolution}.

Untuk setiap $M$ yang merupakan bimodul $(R,R)$, masukkan definisi awal homologi Hochschild
\ref{def:HH} untuk memperoleh
\[
\HHm_n(M)=\Hm_n\left(M\dotimes{R^e}\mathsf{B}R\right)
=\Tor_n^{R^e}(M,R).
\]
\end{contoh}

% segment-id: o014.li2.chapter3.ext-tor.p015
Sekarang kita dapat menggabungkan Contoh \ref{eg:HH-Ext} dan
\ref{eg:HH-Tor} untuk menghitung contoh-contoh homologi dan kohomologi
Hochschild lainnya. Pertama ambil aljabar polinomial $R=\Bbbk[t]$ dan
identifikasi $R^e=R\otimes R^{\opp}$ dengan $\Bbbk[x,y]$. Aksi skalar pada
$R$ sebagai modul-$R^e$ adalah $h(x,y)\cdot f(t)=h(t,t)f(t)$, dan $R$
mempunyai resolusi bebas
% segment-id: o014.li2.chapter3.ext-tor.q007
\[
0\to R^e\xrightarrow{\text{perkalian dengan }x-y}R^e
\xrightarrow{f\mapsto f(t,t)}R\to0.
\]
Dengan menghitung $\Ext^n_{R^e}(R,\cdot)$ dan $\Tor_n^{R^e}(\cdot,R)$ dari
resolusi ini, segera diperoleh $\HHm_n(R)\simeq R\simeq\HHm^n(R)$ untuk
$n\in\{0,1\}$,
sedangkan untuk $n>1$, $\HHm_n(R)=0=\HHm^n(R)$.

% segment-id: o014.li2.chapter3.ext-tor.p016
Selanjutnya ambil $R=\Bbbk[t]/(t^2)$, sehingga
$R^e=\Bbbk[x,y]/(x^2,y^2)$. Untuk modul-$R^e$ $R$ terdapat resolusi bebas
berperiode $2$,
% segment-id: o014.li2.chapter3.ext-tor.q008
\[
\cdots\xrightarrow{x+y}R^e\xrightarrow{x-y}R^e\xrightarrow{x+y}R^e
\xrightarrow{\text{perkalian dengan }x-y}R^e\xrightarrow{f\mapsto f(t,t)}R\to0,
\]
yang keeksakannya kami serahkan kepada pembaca untuk diverifikasi langsung.
Untuk setiap modul-$R$ $M$, mengambil $M\dotimes{R^e}(\cdot)$ dan
$\Hom_{R^e}(\cdot,M)$ masing-masing menghasilkan
% segment-id: o014.li2.chapter3.ext-tor.q009
\begin{gather*}
\cdots\xrightarrow{2t}M\xrightarrow{0}M\xrightarrow{2t}M\xrightarrow{0}M
\to M\dotimes{R^e}R\to0,\\
0\to\Hom_{R^e}(R,M)\to M\xrightarrow{0}M\xrightarrow{2t}M
\xrightarrow{0}M\xrightarrow{2t}\cdots.
\end{gather*}
Maka $\HHm_k(M)$ dan $\HHm^k(M)$ hanya bergantung pada $k\bmod2$;
seandainya hanya dihitung dari definisi, periodisitas ini mungkin sulit
terlihat.

Jika $\Bbbk$ merupakan lapangan, kita dapat membaca $\HHm_n(M)$ dan
$\HHm^n(M)$ dari sini, tetapi jawabannya bergantung pada
$\mathrm{char}(\Bbbk)$. Latihan akan membahas kasus umum
$R=\Bbbk[t]/(t^n)$.

% segment-id: o014.li2.chapter3.ext-tor.p017
Hasil berikut sering digunakan dalam topologi aljabar dan sekaligus menjadi
verifikasi kecil atas materi sebelumnya; hasil ini melibatkan kompleks rantai
beserta homologinya (Catatan \ref{rem:cochain-vs-chain}). Misalkan
$C=(C_\bullet,d^C_\bullet)$ (atau $D=(D_\bullet,d^D_\bullet)$) merupakan
kompleks rantai dari modul-$R$ kanan (atau modul-$R$ kiri). Bentuk
bikompleks rantai $C_\bullet\otimes D_\bullet$ dan tuliskan kompleks totalnya
sebagai
% segment-id: o014.li2.chapter3.ext-tor.q010
\[
C\otimes D:=\tot_{\oplus}\left(C_\bullet\otimes D_\bullet\right);
\]
subskrip $R$ pada hasil kali tensor dihilangkan. Secara konkret, suku
berderajat $n$ dari kompleks total $C\otimes D$ ialah
$\bigoplus_{p+q=n}C_p\otimes D_q$, sedangkan $d_n$ yang dibatasi pada
$C_p\otimes D_q$ sama dengan $d^C_p\otimes\identity+(-1)^p\identity\otimes
d^D_q$. Ini memberikan homomorfisme kanonik yang jelas
\[
\kappa:\bigoplus_{p+q=n}\Hm_p(C)\otimes\Hm_q(D)\to
\Hm_n\left(C\otimes D\right).
\]

% segment-id: o014.li2.chapter3.ext-tor.th002
\begin{teorema}[Teorema Künneth untuk homologi]\label{prop:Kunneth-homology}
\index{kunneth@Teorema Künneth (Künneth Theorem)}
Ambil kompleks rantai $C,D$ seperti di atas. Jika setiap $C_p$ dan
$\Image(d^C_p)$ merupakan modul-$R$ kanan yang datar, maka terdapat barisan
eksak pendek kanonik
% segment-id: o014.li2.chapter3.ext-tor.g001
\[
\begin{tikzcd}[column sep=small, every cell/.append style={font = \small}]
0 \arrow[r] & \displaystyle\bigoplus_{p+q=n} \Hm_p(C) \otimes \Hm_q(D)
\arrow[r, "\kappa"] & \Hm_n\left(C \otimes D \right) \arrow[r] &
\displaystyle\bigoplus_{p+q=n-1} \Tor_1^R\left( \Hm_p(C), \Hm_q(D) \right)
\arrow[r] & 0.
\end{tikzcd}
\]
\end{teorema}

% segment-id: o014.li2.chapter3.ext-tor.pf005
\begin{bukti}
Untuk semua $p\in\Z$, definisikan submodul-submodul dari $C_p$ dengan
$B_p:=\Image(d^C_{p+1})$ dan $Z_p:=\Ker(d^C_p)$. Jadikan keduanya kompleks
rantai $Z=Z_\bullet$ dan $B=B_\bullet$ dengan menetapkan $d^Z=d^B:=0$.
Maka terdapat barisan eksak pendek kompleks rantai
% segment-id: o014.li2.chapter3.ext-tor.g002
\[
\begin{tikzcd}[column sep=large]
0 \arrow[r] & Z \arrow[r] & C \arrow[r, "{d^C=(d^C_p)_p}" inner sep=0.6em]
& B[-1] \arrow[r] & 0.
\end{tikzcd}
\]

Barisan eksak pendek $0\to Z_p\to C_p\xrightarrow{d^C_p}B_{p-1}\to0$
dan barisan eksak panjang funktor $\Tor$ menyiratkan bahwa $Z_p$ juga datar.
Karena $B_{p-1}$ datar, barisan eksak
% segment-id: o014.li2.chapter3.ext-tor.q011
\[
0\to Z_p\otimes D_q\to C_p\otimes D_q\to B_{p-1}\otimes D_q\to0
\]
tetap eksak. Karena jumlah langsung modul (termasuk jumlah tak hingga)
mempertahankan keeksakan, diperoleh barisan eksak pendek kompleks rantai
$0\to Z\otimes D\to C\otimes D\xrightarrow{d^C\otimes\identity_D}
B[-1]\otimes D\to0$.

Karena $d^Z=d^{B[-1]}=0$, barisan eksak panjang homologi yang bersesuaian,
berdasarkan kedataran, menyederhana menjadi
% segment-id: o014.li2.chapter3.ext-tor.q012
\begin{multline*}
\overbracket{\displaystyle\bigoplus_{p+q=n+1}B_{p-1}\otimes\Hm_q(D)}^{=
\bigoplus_{p+q=n}B_p\otimes\Hm_q(D)}\xrightarrow{\partial_{n+1}}
\displaystyle\bigoplus_{p+q=n}Z_p\otimes\Hm_q(D)\to\Hm_n(C\otimes D)\\
\to\underbracket{\displaystyle\bigoplus_{p+q=n}B_{p-1}\otimes\Hm_q(D)}_{=
\bigoplus_{p+q=n-1}B_p\otimes\Hm_q(D)}\xrightarrow{\partial_n}
\displaystyle\bigoplus_{p+q=n-1}Z_p\otimes\Hm_q(D).
\end{multline*}
Yang perlu dijelaskan hanya morfisme penghubung $\partial_{n+1}$ dan
$\partial_n$. Keduanya berasal dari inklusi $B_p\hookrightarrow Z_p$, hingga
kemungkinan beberapa tanda positif atau negatif. Untuk melihatnya, kembali ke
konstruksi eksplisit barisan eksak panjang, yaitu versi homologi dari Proposisi
\ref{prop:long-exact-sequence-ses}; di sini kita dapat menggunakan penelusuran
diagram; lihat uraian sebelum \cite[Proposisi 6.8.6]{Li1}. Detailnya diserahkan
kepada pembaca.

Sebagai korolari, $\Coker(\partial_{n+1})$ dapat diidentifikasi dengan
$\bigoplus_{p+q=n}\Hm_p(C)\otimes\Hm_q(D)$; monomorfisme darinya ke
$\Hm_n(C\otimes D)$ adalah $\kappa$.

Akhirnya, $0\to B_p\to Z_p\to\Hm_p(C)\to0$ merupakan resolusi datar
$\Hm_p(C)$. Karena itu,
$\Ker(\partial_n)\simeq\bigoplus_{p+q=n-1}\Tor_1^R(\Hm_p(C),\Hm_q(D))$.
Ini menghasilkan barisan eksak pendek yang dinyatakan.
\end{bukti}

% segment-id: o014.li2.chapter3.ext-tor.p018
Perhatikan bahwa jika $R$ merupakan gelanggang pembagian, maka $\kappa$ pasti
merupakan isomorfisme.

% segment-id: o014.li2.chapter3.ext-tor.p019
Karena $D$ berperan sebagai koefisien homologi dalam konteks topologi, Teorema
\ref{prop:Kunneth-homology} pada kasus khusus ketika $D$ terkonsentrasi pada
suku berderajat nol juga disebut \emph{teorema koefisien universal}.
Hasil-hasil semacam ini memiliki berbagai variasi; selain itu, struktur
gelanggang memberikan operasi pada $\Tor$ yang tidak dimiliki bifunktor
turunan lain. Struktur tersebut tampaknya lebih tepat dirapikan dalam bahasa
kategori turunan atau barisan spektral.
\index{teorema koefisien universal@teorema koefisien universal (Universal Coefficient Theorem)}
